\documentclass[12pt]{article}

% Idioma y codificación
\usepackage[spanish, es-tabla, es-notilde]{babel}       %es-tabla para que se titule "Tabla"
\usepackage[utf8]{inputenc}

% Márgenes
\usepackage[a4paper,top=3cm,bottom=2.5cm,left=3cm,right=3cm]{geometry}

% Comentarios de bloque
\usepackage{verbatim}

% Paquetes de links
\usepackage[hidelinks]{hyperref}    % Permite enlaces
\usepackage{url}                    % redirecciona a la web

% Más opciones para enumeraciones
\usepackage{enumitem}

% Personalizar la portada
\usepackage{titling}

% Paquetes de tablas
\usepackage{multirow}

% Para añadir el símbolo de euro
\usepackage{eurosym}

% Paquete para subrayar con diferentes estilos
\usepackage[normalem]{ulem}


%------------------------------------------------------------------------

%Paquetes de figuras
\usepackage{caption}
\usepackage{subcaption} % Figuras al lado de otras
\usepackage{float}      % Poner figuras en el sitio indicado H.


% Paquetes de imágenes
\usepackage{graphicx}       % Paquete para añadir imágenes
\usepackage{transparent}    % Para manejar la opacidad de las figuras

% Paquete para usar colores
\usepackage[dvipsnames, table, xcdraw]{xcolor}
\usepackage{pagecolor}      % Para cambiar el color de la página

% Habilita tamaños de fuente mayores
\usepackage{fix-cm}

% Para los gráficos
\usepackage{tikz}
\usepackage{forest}

% Para poder situar los nodos en los grafos
\usetikzlibrary{positioning}

% Para evitar problemas entre babel y tikz
\usetikzlibrary{babel}


%------------------------------------------------------------------------

% Paquetes de matemáticas
\usepackage{mathtools, amsfonts, amssymb, mathrsfs}
\usepackage[makeroom]{cancel}     % Simplificar tachando
\usepackage{polynom}    % Divisiones y Ruffini
\usepackage{units} % Para poner fracciones diagonales con \nicefrac

\usepackage{pgfplots}   %Representar funciones
\pgfplotsset{compat=1.18}  % Versión 1.18

\usepackage{tikz-cd}    % Para usar diagramas de composiciones
\usetikzlibrary{calc}   % Para usar cálculo de coordenadas en tikz

%Definición de teoremas, etc.
\usepackage{amsthm}
%\swapnumbers   % Intercambia la posición del texto y de la numeración

\theoremstyle{plain}

\makeatletter
\@ifclassloaded{article}{
  % --- Configuración para ARTICLE ---
  \newtheorem{teo}{Teorema}[section]
  \theoremstyle{definition}
  \newtheorem{definicion}{Definición}[section]
  \newtheorem{ejercicio}{Ejercicio}[section]

  \renewcommand{\theteo}{%
    \ifnum\value{section}=0 \arabic{teo}\else \thesection.\arabic{teo}\fi
  }
  \renewcommand{\thedefinicion}{%
    \ifnum\value{section}=0 \arabic{definicion}\else \thesection.\arabic{definicion}\fi
  }
  \renewcommand{\theejercicio}{%
    \ifnum\value{section}=0 \arabic{ejercicio}\else \thesection.\arabic{ejercicio}\fi
  }
}{
  % --- Configuración para BOOK / REPORT ---
  \newtheorem{teo}{Teorema}[chapter]
  \theoremstyle{definition}
  \newtheorem{definicion}{Definición}[chapter]
  \newtheorem{ejercicio}{Ejercicio}[section]

  % Si no hay sección: Capítulo.Teorema (ej. 2.1)
  % Si hay sección:    Capítulo.Sección.Teorema (ej. 2.3.1)
  \renewcommand{\theteo}{%
    \ifnum\value{section}=0
      \thechapter.\arabic{teo}%
    \else
      \thesection.\arabic{teo}% (\thesection ya contiene el número de capítulo)
    \fi
  }
  \renewcommand{\thedefinicion}{%
    \ifnum\value{section}=0
      \thechapter.\arabic{definicion}%
    \else
      \thesection.\arabic{definicion}%
    \fi
  }
  \renewcommand{\theejercicio}{%
    \ifnum\value{section}=0
      \arabic{ejercicio}%
    \else
      \thesection.\arabic{ejercicio}% Formato por defecto: ej. 2.2.1
    \fi
  }
}
\makeatother

% Teoremas dependientes o con contador compartido
\newtheorem{coro}{Corolario}[teo]
\newtheorem{prop}[teo]{Proposición}
\newtheorem{lema}[teo]{Lema}
\newtheorem*{lema*}{Lema}

\theoremstyle{remark}
\newtheorem*{observacion}{Observación}

\theoremstyle{definition}
\newtheorem*{notacion}{Notación}
\newtheorem*{ejemplo}{Ejemplo}
\newtheorem*{ejercicio*}{Ejercicio}


% \renewcommand\qedsymbol{$\blacksquare$}         % Cambiar símbolo QED
%------------------------------------------------------------------------

% Paquetes para encabezados
\usepackage{fancyhdr}
\pagestyle{fancy}
\fancyhf{}

\newcommand{\helv}{ % Modificación tamaño de letra
\fontfamily{}\fontsize{12}{12}\selectfont}
\setlength{\headheight}{15pt} % Amplía el tamaño del índice


%\usepackage{lastpage}   % Referenciar última pag   \pageref{LastPage}
%\fancyfoot[C]{%
%  \begin{minipage}{\textwidth}
%    \centering
%    ~\\
%    \thepage\\
%    \href{https://losdeldgiim.github.io/}{\texttt{\footnotesize losdeldgiim.github.io}}
%  \end{minipage}
%}
\fancyfoot[C]{\thepage}
\fancyfoot[R]{\href{https://losdeldgiim.github.io/}{\texttt{\footnotesize losdeldgiim.github.io}}}

%------------------------------------------------------------------------

% Conseguir que no ponga "Capítulo 1". Sino solo "1."
\makeatletter
\@ifclassloaded{book}{
  \renewcommand{\chaptermark}[1]{\markboth{\thechapter.\ #1}{}} % En el encabezado
    
  \renewcommand{\@makechapterhead}[1]{%
  \vspace*{50\p@}%
  {\parindent \z@ \raggedright \normalfont
    \ifnum \c@secnumdepth >\m@ne
      \huge\bfseries \thechapter.\hspace{1em}\ignorespaces
    \fi
    \interlinepenalty\@M
    \Huge \bfseries #1\par\nobreak
    \vskip 40\p@
  }}
}
\makeatother

%------------------------------------------------------------------------
% Paquetes de cógido
\usepackage{minted}
\renewcommand\listingscaption{Código fuente}

\usepackage{fancyvrb}
% Personaliza el tamaño de los números de línea
\renewcommand{\theFancyVerbLine}{\small\arabic{FancyVerbLine}}

% Estilo para C++
\newminted{cpp}{
    frame=lines,
    framesep=2mm,
    baselinestretch=1.2,
    linenos,
    escapeinside=||
}

% para minted
\definecolor{LightGray}{rgb}{0.95,0.95,0.92}
\setminted{
    linenos=true,
    stepnumber=5,
    numberfirstline=true,
    autogobble,
    breaklines=true,
    breakautoindent=true,
    breaksymbolleft=,
    breaksymbolright=,
    breaksymbolindentleft=0pt,
    breaksymbolindentright=0pt,
    breaksymbolsepleft=0pt,
    breaksymbolsepright=0pt,
    fontsize=\footnotesize,
    bgcolor=LightGray,
    numbersep=10pt
}


\usepackage{listings} % Para incluir código desde un archivo

\renewcommand\lstlistingname{Código Fuente}
\renewcommand\lstlistlistingname{Índice de Códigos Fuente}

% Definir colores
\definecolor{vscodepurple}{rgb}{0.5,0,0.5}
\definecolor{vscodeblue}{rgb}{0,0,0.8}
\definecolor{vscodegreen}{rgb}{0,0.5,0}
\definecolor{vscodegray}{rgb}{0.5,0.5,0.5}
\definecolor{vscodebackground}{rgb}{0.97,0.97,0.97}
\definecolor{vscodelightgray}{rgb}{0.9,0.9,0.9}

% Configuración para el estilo de C similar a VSCode
\lstdefinestyle{vscode_C}{
  backgroundcolor=\color{vscodebackground},
  commentstyle=\color{vscodegreen},
  keywordstyle=\color{vscodeblue},
  numberstyle=\tiny\color{vscodegray},
  stringstyle=\color{vscodepurple},
  basicstyle=\scriptsize\ttfamily,
  breakatwhitespace=false,
  breaklines=true,
  captionpos=b,
  keepspaces=true,
  numbers=left,
  numbersep=5pt,
  showspaces=false,
  showstringspaces=false,
  showtabs=false,
  tabsize=2,
  frame=tb,
  framerule=0pt,
  aboveskip=10pt,
  belowskip=10pt,
  xleftmargin=10pt,
  xrightmargin=10pt,
  framexleftmargin=10pt,
  framexrightmargin=10pt,
  framesep=0pt,
  rulecolor=\color{vscodelightgray},
  backgroundcolor=\color{vscodebackground},
}

%------------------------------------------------------------------------
% La etiqueda IA sirve para indicar a partir de donde le pasamos el preambulo a la IA para el escaneo
% IA

% Comandos definidos
\newcommand{\bb}[1]{\mathbb{#1}}
\newcommand{\cc}[1]{\mathcal{#1}}

% I prefer the slanted \leq
\let\oldleq\leq % save them in case they're every wanted
\let\oldgeq\geq
\renewcommand{\leq}{\leqslant}
\renewcommand{\geq}{\geqslant}

% Si y solo si
\newcommand{\sii}{\iff}

% MCD y MCM
\DeclareMathOperator{\mcd}{mcd}
\DeclareMathOperator{\mcm}{mcm}

% Signo
\DeclareMathOperator{\sgn}{sgn}

% Imagen
\DeclareMathOperator{\Ima}{Im}

% Aplicación identidad
\newcommand{\id}{\mathrm{Id}}

% Letras griegas
\newcommand{\eps}{\epsilon}
\newcommand{\veps}{\varepsilon}
\newcommand{\lm}{\lambda}

\newcommand{\ol}{\overline}
\newcommand{\ul}{\underline}
\newcommand{\wt}{\widetilde}
\newcommand{\wh}{\widehat}

\let\oldvec\vec
\renewcommand{\vec}{\overrightarrow}

% Derivadas parciales
\newcommand{\del}[2]{\frac{\partial #1}{\partial #2}}
\newcommand{\Del}[3]{\frac{\partial^{#1} #2}{\partial #3^{#1}}}
\newcommand{\deld}[2]{\dfrac{\partial #1}{\partial #2}}
\newcommand{\Deld}[3]{\dfrac{\partial^{#1} #2}{\partial #3^{#1}}}


\newcommand{\AstIg}{\stackrel{(\ast)}{=}}
\newcommand{\Hop}{\stackrel{L'H\hat{o}pital}{=}}

\newcommand{\red}[1]{{\color{red}#1}} % Para integrales, destacar los cambios.

% Método de integración
\newcommand{\MetInt}[2]{
    \left[\begin{array}{c}
        #1 \\ #2
    \end{array}\right]
}

% Declarar aplicaciones
% 1. Nombre aplicación
% 2. Dominio
% 3. Codominio
% 4. Variable
% 5. Imagen de la variable
\newcommand{\Func}[5]{
    \begin{equation*}
        \begin{array}{rrll}
            \displaystyle #1:& \displaystyle  #2 & \longrightarrow & \displaystyle  #3\\
               & \displaystyle  #4 & \longmapsto & \displaystyle  #5
        \end{array}
    \end{equation*}
}

%------------------------------------------------------------------------


\newcommand{\myhref}[2]{\href{#1}{\color{blue}\uline{#2}}}


\begin{document}

    % 1. Foto de fondo
    % 2. Título
    % 3. Encabezado Izquierdo
    % 4. Color de fondo
    % 5. Coord x del titulo
    % 6. Coord y del titulo
    % 7. Fecha

    
    % 1. Foto de fondo
% 2. Título
% 3. Encabezado Izquierdo
% 4. Color de fondo
% 5. Coord x del titulo
% 6. Coord y del titulo
% 7. Fecha
% 8. Autor

\newcommand{\portada}[8]{

    \portadaBase{#1}{#2}{#3}{#4}{#5}{#6}{#7}{#8}
    \portadaBook{#1}{#2}{#3}{#4}{#5}{#6}{#7}{#8}
}

\newcommand{\portadaExamen}[8]{

    \portadaBase{#1}{#2}{#3}{#4}{#5}{#6}{#7}{#8}
    \portadaArticle{#1}{#2}{#3}{#4}{#5}{#6}{#7}{#8}
}




\newcommand{\portadaBase}[8]{

    % Tiene la portada principal y la licencia Creative Commons
    
    % 1. Foto de fondo
    % 2. Título
    % 3. Encabezado Izquierdo
    % 4. Color de fondo
    % 5. Coord x del titulo
    % 6. Coord y del titulo
    % 7. Fecha
    % 8. Autor
    
    
    \thispagestyle{empty}               % Sin encabezado ni pie de página
    \newgeometry{margin=0cm}        % Márgenes nulos para la primera página
    
    
    % Encabezado
    \fancyhead[L]{\helv #3}
    \fancyhead[R]{\helv \nouppercase{\leftmark}}
    
    
    \pagecolor{#4}        % Color de fondo para la portada
    
    \begin{figure}[p]
        \centering
        \transparent{0.3}           % Opacidad del 30% para la imagen
        
        \includegraphics[width=\paperwidth, keepaspectratio]{../../_assets/#1}
    
        \begin{tikzpicture}[remember picture, overlay]
            \node[anchor=north west, text=white, opacity=1, font=\fontsize{60}{90}\selectfont\bfseries\sffamily, align=left] at (#5, #6) {#2};
            
            \node[anchor=south east, text=white, opacity=1, font=\fontsize{12}{18}\selectfont\sffamily, align=right] at (9.7, 3) {\href{https://losdeldgiim.github.io/}{\textbf{Los Del DGIIM}, \texttt{\footnotesize losdeldgiim.github.io}}};
            
            \node[anchor=south east, text=white, opacity=1, font=\fontsize{12}{15}\selectfont\sffamily, align=right] at (9.7, 1.8) {Doble Grado en Ingeniería Informática y Matemáticas\\Universidad de Granada};
        \end{tikzpicture}
    \end{figure}
    
    
    \restoregeometry        % Restaurar márgenes normales para las páginas subsiguientes
    \pagecolor{white}       % Restaurar el color de página
    
    
    \newpage
    \thispagestyle{empty}               % Sin encabezado ni pie de página
    \begin{tikzpicture}[remember picture, overlay]
        \node[anchor=south west, inner sep=3cm] at (current page.south west) {
            \begin{minipage}{0.5\paperwidth}
                \href{https://creativecommons.org/licenses/by-nc-nd/4.0/}{
                    \includegraphics[height=2cm]{../../_assets/Licencia.png}
                }\vspace{1cm}\\
                Esta obra está bajo una
                \href{https://creativecommons.org/licenses/by-nc-nd/4.0/}{
                    Licencia Creative Commons Atribución-NoComercial-SinDerivadas 4.0 Internacional (CC BY-NC-ND 4.0).
                }\\
    
                Eres libre de compartir y redistribuir el contenido de esta obra en cualquier medio o formato, siempre y cuando des el crédito adecuado a los autores originales y no persigas fines comerciales. 
            \end{minipage}
        };
    \end{tikzpicture}
    
    
    
    % 1. Foto de fondo
    % 2. Título
    % 3. Encabezado Izquierdo
    % 4. Color de fondo
    % 5. Coord x del titulo
    % 6. Coord y del titulo
    % 7. Fecha
    % 8. Autor


}


\newcommand{\portadaBook}[8]{

    % 1. Foto de fondo
    % 2. Título
    % 3. Encabezado Izquierdo
    % 4. Color de fondo
    % 5. Coord x del titulo
    % 6. Coord y del titulo
    % 7. Fecha
    % 8. Autor

    % Personaliza el formato del título
    \pretitle{\begin{center}\bfseries\fontsize{42}{56}\selectfont}
    \posttitle{\par\end{center}\vspace{2em}}
    
    % Personaliza el formato del autor
    \preauthor{\begin{center}\Large}
    \postauthor{\par\end{center}\vfill}
    
    % Personaliza el formato de la fecha
    \predate{\begin{center}\huge}
    \postdate{\par\end{center}\vspace{2em}}
    
    \title{#2}
    \author{\href{https://losdeldgiim.github.io/}{Los Del DGIIM, \texttt{\large losdeldgiim.github.io}}
    \\ \vspace{0.5cm}#8}
    \date{Granada, #7}
    \maketitle
    
    \tableofcontents
}




\newcommand{\portadaArticle}[8]{

    % 1. Foto de fondo
    % 2. Título
    % 3. Encabezado Izquierdo
    % 4. Color de fondo
    % 5. Coord x del titulo
    % 6. Coord y del titulo
    % 7. Fecha
    % 8. Autor

    % Personaliza el formato del título
    \pretitle{\begin{center}\bfseries\fontsize{42}{56}\selectfont}
    \posttitle{\par\end{center}\vspace{2em}}
    
    % Personaliza el formato del autor
    \preauthor{\begin{center}\Large}
    \postauthor{\par\end{center}\vspace{3em}}
    
    % Personaliza el formato de la fecha
    \predate{\begin{center}\huge}
    \postdate{\par\end{center}\vspace{5em}}
    
    \title{#2}
    \author{\href{https://losdeldgiim.github.io/}{Los Del DGIIM, \texttt{\large losdeldgiim.github.io}}
    \\ \vspace{0.5cm}#8}
    \date{Granada, #7}
    \thispagestyle{empty}               % Sin encabezado ni pie de página
    \maketitle
    \vfill
}
    \portadaExamen{ffccA4.jpg}{Análisis Funcional\\Examen V}{Análisis Funcional. Examen V}{MidnightBlue}{-8}{28}{2025}{}

    \begin{description}
        \item[Asignatura] Análisis Funcional.
        \item[Curso Académico] 2021/22.
        \item[Grado] Doble Grado en Ingeniería Informática y Matemáticas.
        % \item[Grupo] ---.
        % \item[Profesor] ---.
        \item[Descripción] Examen Ordinario.
        % \item[Fecha] 15 de enero de 2025.
        % \item[Duración] ---.
    
    \end{description}
    \newpage


    % ------------------------------------
    
    \begin{ejercicio}[1.5 puntos]
        Sea $X$ el espacio normado de las funciones continuas en $[0,1]$ que se anulan en cero con la norma del máximo:
        \begin{equation*}
            X = \{f\in C[0,1] : f(0) = 0\}, \qquad \|f\|_\infty = \max\{|f(t)| : 0\leq t \leq 1\}
        \end{equation*}
        En $X$ se considera el funcional lineal $\varphi:X\to \mathbb{R}$ dado por:
        \begin{equation*}
            \varphi(f) = \int_{0}^{1} tf(t)~dt \qquad \forall f\in X
        \end{equation*}
        \begin{enumerate}[label=\alph*)]
            \item \textbf{[1 punto]} Calcula $\|\varphi\|$.
            \item \textbf{[0.5 puntos]} Prueba que $\varphi$ no alcanza su norma.
        \end{enumerate}
    \end{ejercicio}

    \begin{ejercicio}[2 puntos]
        Sea $T:c_0\to l_2$ el operador lineal definido para todo $x\in c_0$ por
        \begin{equation*}
            [Tx](n) = \frac{x(n)}{n}\qquad \forall n\in \mathbb{N}
        \end{equation*}
        \begin{enumerate}[label=\alph*)]
            \item \textbf{[1 punto]} Calcula su norma y prueba que $T$ no la alcanza.
            \item \textbf{[0.5 punos]} Prueba que $Y=T(c_0)$ es un subespacio denso en $l_2$ y que $T$ es una biyección lineal de $c_0$ sobre $Y$.
            \item \textbf{[0.5 puntos]} ¿Es continua la aplicación $T^{-1}:Y\to c_0$? ¿Es $Y$ cerrado en $l_2$?
        \end{enumerate}
    \end{ejercicio}

    \begin{ejercicio}[1.5 puntos]
        Sean $X$ un espacio de Banach, $Y$ un espacio normado y $\cc{F}$ un subconjunto de $L(X,Y)$. Prueba que las siguientes afirmaciones son equivalentes:
        \begin{enumerate}[label=\alph*)]
            \item $\cc{F}$ está acotado.
            \item Para cada $x\in X$ y cada $g\in Y^\ast$ el conjunto $\{g(T(x)) : T\in \cc{F}\}$ está acotado.
        \end{enumerate}
    \end{ejercicio}

    \begin{ejercicio}[2 puntos]
        Responde a las siguientes cuestiones indicando en cada caso, según proceda, un resultado de teoría que justifique tu respuesta o un contraejemplo, o bien dando una prueba.
        \begin{enumerate}[label=\alph*)]
            \item \textbf{[0.5 puntos]} ¿Puede definirse alguna norma completa en el espacio vectorial $c_{00}$ de las sucesiones casi nulas?
            \item \textbf{[0.5 puntos]} Si $H$ es un espacio de Hilbert y $M$ es un subespacio de $H$ verificando que $M^\perp = \{0\}$, entonces $M = H$.
            \item \textbf{[0.5 puntos]} Sea $\{f_n\}$ una sucesión puntualmente acotada de funciones lineales continuos en un espacio de Banach $X$, definiendo $T(x) = \{f_n(x)\}$ se obtiene un operador lineal continuo de $X$ en $t_\infty$.
            \item \textbf{[0.5 puntos]} Sean $\|\cdot \|$ y $|\|\cdot |\|$ dos normas completas no equivalentes en un espacio vectorial $X$. Entonces la aplicación identidad $I_X:(X,\|\cdot \|)\to (X,|\|\cdot |\|)$ no es continua.
        \end{enumerate}
    \end{ejercicio}

    \begin{ejercicio}[3 puntos]
        Escribe sobre alguno de los siguientes temas:
        \begin{enumerate}[label=\alph*)]
            \item Ortogonalidad. Teorema de la proyección ortogonal. Teorema de Riesz-Fréchet.
            \item Separación de conjuntos convexos en espacios normados.
            \item Lema de categoría de Baire. Teorema de Banach-Steinhaus.
        \end{enumerate}
    \end{ejercicio}

    \newpage
    \setcounter{ejercicio}{0} % Reiniciar contador de ejercicios

    \begin{ejercicio}[1.5 puntos]
        Sea $X$ el espacio normado de las funciones continuas en $[0,1]$ que se anulan en cero con la norma del máximo:
        \begin{equation*}
            X = \{f\in C[0,1] : f(0) = 0\}, \qquad \|f\|_\infty = \max\{|f(t)| : 0\leq t \leq 1\}
        \end{equation*}
        En $X$ se considera el funcional lineal $\varphi:X\to \mathbb{R}$ dado por:
        \begin{equation*}
            \varphi(f) = \int_{0}^{1} tf(t)~dt \qquad \forall f\in X
        \end{equation*}
        \begin{enumerate}[label=\alph*)]
            \item \textbf{[1 punto]} Calcula $\|\varphi\|$.
            
            Veamos en primer lugar que dicho funcional es continuo. Para ello, sea $f\in X$. Entonces:
            \begin{equation*}
                |\varphi(f)| = \left|\int_0^1 tf(t)~dt\right| \leq \int_0^1 |tf(t)|~dt \leq \int_0^1 t\|f\|_\infty~dt = \|f\|_\infty \int_0^1 t~dt = \frac{1}{2}\|f\|_\infty
            \end{equation*}
            Por lo tanto, $\|\varphi\| \leq \nicefrac{1}{2}$. Para ver que $\|\varphi\| = \nicefrac{1}{2}$, consideremos la sucesión de funciones $\{f_n\}$ mostrada en la Figura~\ref{fig:funciones}. Cada función $f_n$ está definida por:
            \begin{equation*}
                f_n(t) = \begin{cases}
                    nt & \text{si } t\in \left[0, \nicefrac{1}{n}\right], \\
                    1 & \text{si } t\in \left[\nicefrac{1}{n}, 1\right]
                \end{cases}
            \end{equation*}
            \begin{figure}[h]
                \centering
                \begin{tikzpicture}
                    \begin{axis}[
                        axis lines=middle,
                        xmin=0, xmax=1.1,
                        ymin=0, ymax=1.2,
                        xtick={0, 0.5, 1},
                        xticklabels={$0$, $\nicefrac{1}{2}$, $1$},
                        ytick={0, 1},
                        yticklabels={$0$, $1$},
                        width=10cm,
                        height=6cm,
                        legend pos=south east
                    ]
                        \foreach \n/\col in {2/blue!70, 4/red!70, 8/green!60!black} {
                            \edef\temp{
                                \noexpand\addplot[thick, color=\col] coordinates {
                                    (0, 0)
                                    (1/\n, 1)
                                    (1, 1)
                                };
                                \noexpand\addlegendentry{$n=\n$}

                                % \draw[dashed, color=\col] (axis cs:1/\n,0) -- (axis cs:1/\n,1);
                            }
                            \temp
                        }
                        \draw[dashed, gray] (axis cs:0,1) -- (axis cs:1,1);
                    \end{axis}
                \end{tikzpicture}
                \caption{Funciones $f_n$ para $n=2,4,8$.}
                \label{fig:funciones}                
            \end{figure}

            Se verifica que $\|f_n\|_\infty = 1$ para todo $n$. Además, se tiene:
            \begin{align*}
                \varphi(f_n) &= \int_0^{\nicefrac{1}{n}} t(nt)~dt + \int_{\nicefrac{1}{n}}^1 t(1)~dt = \int_0^{\nicefrac{1}{n}} n t^2~dt + \int_{\nicefrac{1}{n}}^1 t~dt = n\left[\frac{t^3}{3}\right]_0^{\nicefrac{1}{n}} + \left[\frac{t^2}{2}\right]_{\nicefrac{1}{n}}^1 =\\&= \frac{1}{3n^2} + \left(\frac{1}{2} - \frac{1}{2n^2}\right) = \frac{1}{2} - \frac{1}{6n^2}.
            \end{align*}

            Por lo tanto:
            \begin{equation*}
                \|\varphi\| \geq \dfrac{|\varphi(f_n)|}{\|f_n\|_\infty} = \frac{1}{2} - \frac{1}{6n^2}\qquad \forall n\geq 2
            \end{equation*}
            Tomando el límite cuando $n\to\infty$, obtenemos que $\|\varphi\| \geq \nicefrac{1}{2}$. Combinando esto con la desigualdad anterior, concluimos que $\|\varphi\| = \nicefrac{1}{2}$.
            \item \textbf{[0.5 puntos]} Prueba que $\varphi$ no alcanza su norma.\\
            
            Supongamos que existe $f\in X$ tal que $\|f\|_\infty = 1$ y $|\varphi(f)| = \nicefrac{1}{2}$. Entonces:
            \begin{equation*}
                \frac{1}{2} = |\varphi(f)| = \left|\int_0^1 tf(t)~dt\right| \leq \int_0^1 |tf(t)|~dt \leq \int_0^1 t\|f\|_\infty~dt = \frac{1}{2}
            \end{equation*}
            Por lo tanto, todas las desigualdades deben ser iguales. En particular,
            \begin{equation*}
                \int_0^1 |tf(t)|~dt = \int_0^1 t\|f\|_\infty~dt \implies \int_0^1 t\left(\|f\|_\infty - |f(t)|\right)~dt = 0
            \end{equation*}
            Como $t\geq 0$ y $\|f\|_\infty - |f(t)| \geq 0$ para todo $t\in [0,1]$, y haciendo uso de que $f$ es continua, se deduce que:
            \begin{equation*}
                t\left(\|f\|_\infty - |f(t)|\right) = 0 \qquad \forall t\in [0,1]
            \end{equation*}
            Por lo tanto, para todo $t\in \left]0,1\right]$, se tiene que $|f(t)| = \|f\|_\infty = 1$. Puesto que $f$ es continua con $f(0) = 0$, y el módulo de $f$ es una función continua, se tiene que:
            \begin{equation*}
                0 = |f(0)| = \left|\lim_{t\to 0^+} f(t)\right| = \lim_{t\to 0^+} |f(t)| = \lim_{t\to 0^+} 1 = 1
            \end{equation*}

            Esto es una contradicción, por lo que no existe tal función $f$ y $\varphi$ no alcanza su norma.
        \end{enumerate}
    \end{ejercicio}

    \begin{ejercicio}[2 puntos]
        Sea $T:c_0\to l_2$ el operador lineal definido para todo $x\in c_0$ por
        \begin{equation*}
            [Tx](n) = \frac{x(n)}{n}\qquad \forall n\in \mathbb{N}
        \end{equation*}
        \begin{enumerate}[label=\alph*)]
            \item \textbf{[1 punto]} Calcula su norma y prueba que $T$ no la alcanza.
            
            Tenemos que:
            \begin{equation*}
                \|Tx\|_2^2 = \sum_{n=1}^{\infty} \left|\frac{x(n)}{n}\right|^2 \leq \sum_{n=1}^{\infty} \frac{\|x\|_\infty^2}{n^2} = \|x\|_\infty^2 \sum_{n=1}^{\infty} \frac{1}{n^2}
            \end{equation*}

            Como $2>1$, la serie armónica dada es convergente:
            \begin{equation*}
                \sum_{n=1}^{\infty} \frac{1}{n^2} = M<\infty
            \end{equation*}

            De hecho, se sabe que $M = \nicefrac{\pi^2}{6}$, aunque no es necesario para el cálculo de la norma. Por lo tanto:
            \begin{equation*}
                \|Tx\|_2^2 \leq M\|x\|_\infty^2 \implies \|Tx\|_2 \leq \sqrt{M}\|x\|_\infty
            \end{equation*}

            Por lo tanto, $T$ es continuo y $\|T\| \leq \sqrt{M}$. Para ver que $\|T\| = \sqrt{M}$, consideremos la sucesión $\{x_n\}$ en $c_{00}\subset c_0$ dada por:
            \begin{equation*}
                x_n(k) = \begin{cases}
                    1 & \text{si } k\leq n, \\
                    0 & \text{si } k>n
                \end{cases}
            \end{equation*}

            Entonces $\|x_n\|_\infty = 1$ y:
            \begin{equation*}
                \|Tx_n\|_2^2 = \sum_{k=1}^{\infty} \left|\frac{x_n(k)}{k}\right|^2 = \sum_{k=1}^{n} \frac{1}{k^2}
            \end{equation*}

            Por lo tanto:
            \begin{equation*}
                \|T\| \geq \frac{\|Tx_n\|_2}{\|x_n\|_\infty} = \sqrt{\sum_{k=1}^{n} \frac{1}{k^2}} \qquad \forall n\in \mathbb{N}
            \end{equation*}

            Tomando el límite cuando $n\to\infty$, puesto que la raíz cuadrada es continua, obtenemos:
            \begin{equation*}
                \|T\| \geq \sqrt{\sum_{k=1}^{\infty} \frac{1}{k^2}} = \sqrt{M}
            \end{equation*}

            Combinando esto con la desigualdad anterior, concluimos que $\|T\| = \sqrt{M}$.\\

            Veamos ahora que $T$ no alcanza su norma. Supongamos que existe $x\in c_0$ tal que $\|x\|_\infty = 1$ y $\|Tx\|_2 = \sqrt{M}$. Entonces:
            \begin{equation*}
                \sqrt{M} = \|Tx\|_2 = \sqrt{\sum_{n=1}^{\infty} \left|\frac{x(n)}{n}\right|^2} \leq \sqrt{\sum_{n=1}^{\infty} \frac{\|x\|_\infty^2}{n^2}} = \sqrt{M}
            \end{equation*}

            Por tanto, todas las desigualdades deben ser iguales. En particular:
            \begin{equation*}
                \sum_{n=1}^{\infty} \left|\frac{x(n)}{n}\right|^2 = \sum_{n=1}^{\infty} \frac{\|x\|_\infty^2}{n^2} \implies |x(n)| = \|x\|_\infty = 1 \qquad \forall n\in \mathbb{N}
            \end{equation*}

            Esto es una contradicción, ya que $x\in c_0$ y por tanto $\{x(n)\}$ debe tender a cero. Por lo tanto, $T$ no alcanza su norma.

            \item \textbf{[0.5 punos]} Prueba que $Y=T(c_0)$ es un subespacio denso en $l_2$ y que $T$ es una biyección lineal de $c_0$ sobre $Y$.
            
            Veamos en primer lugar que $T$ es una biyección lineal de $c_0$ sobre $Y$. La linealidad de $T$ es inmediata por la definición de $T$. Además, es claro que $T$ es sobreyectiva de $c_0$ sobre $Y=T(c_0)$ por definición de $Y$. Por último, veamos que $T$ es inyectiva. Supongamos que $T(x)=0$ para algún $x\in c_0$. Entonces:
            \begin{equation*}
                T(x)(n) = \frac{x(n)}{n} = 0 \qquad \forall n\in \mathbb{N}
            \end{equation*}
            Por lo tanto, $x(n) = 0$ para todo $n\in \mathbb{N}$, es decir, $x=0$. Por lo tanto, $T$ es inyectiva y por tanto una biyección lineal de $c_0$ sobre $Y$. Veamos ahora que $Y$ es denso en $l_2$. Sea $y\in l_2$, y queremos ver que $y\in \overline{Y}$, lo que equivale a ver que existe una sucesión $\{y_n\}\subset Y=T(c_0)$ tal que $\{y_n\}\to y$ en $l_2$. Consideremos la sucesión $\{x_n\}\subset c_{00}\subset c_0$ dada por:
            \begin{equation*}
                x_n(k) = \begin{cases}
                    k y(k) & \text{si } k\leq n, \\
                    0 & \text{si } k>n
                \end{cases}
            \end{equation*}
            
            Consideramos ahora la sucesión $\{y_n\}\subset Y$ dada por $y_n = T(x_n)$. Entonces:
            \begin{equation*}
                y_n(k) = T(x_n)(k) = \frac{x_n(k)}{k} = \begin{cases}
                    y(k) & \text{si } k\leq n, \\
                    0 & \text{si } k>n
                \end{cases}
            \end{equation*}

            Por lo tanto, $\{y_n\}$ es la sucesión de truncamientos de $y$. Veamos ahora que $\{y_n\}\to y$ en $l_2$. Para ello, consideremos:
            \begin{equation*}
                \|y_n - y\|_2^2 = \sum_{k=1}^{\infty} |y_n(k) - y(k)|^2 = \sum_{k=n+1}^{\infty} |y(k)|^2
            \end{equation*}

            Como $y\in l_2$, se tiene que $\sum\limits_{k=1}^{\infty} |y(k)|^2 < \infty$, y por tanto:
            \begin{equation*}
                \lim_{n\to\infty} \|y_n - y\|_2^2 = \lim_{n\to\infty} \sum_{k=n+1}^{\infty} |y(k)|^2 = 0
            \end{equation*}

            Esto prueba que $\{y_n\}\to y$ en $l_2$, y por tanto $y\in \overline{Y}$. Como $y\in l_2$ era arbitrario, se concluye que $\overline{Y} = l_2$, es decir, $Y$ es denso en $l_2$.

            \item \textbf{[0.5 puntos]} ¿Es continua la aplicación $T^{-1}:Y\to c_0$? ¿Es $Y$ cerrado en $l_2$?
            
            Veamos que $T^{-1}$ no es continua. Esta aplicación está definida por:
            \begin{equation*}
                T^{-1}(y)(n) = n y(n) \qquad \forall y\in Y, \forall n\in \mathbb{N}
            \end{equation*}

            Consideremos la sucesión $\{ne_n\}\subset c_{00}\subset c_0$, donde $e_n$ es la sucesión canónica de $l_2$. Entonces, $T(ne_n) = e_n \in Y$, luego consideramos la sucesión $\{e_n\}\subset Y$.
            \begin{align*}
                \|T^{-1}(e_n)\|_\infty = \|ne_n\|_\infty = n
            \end{align*}

            Por otro lado, $\|e_n\|_2 = 1$ para todo $n$. Por lo tanto, $\{e_n\}$ es una sucesión acotada en $Y$, pero $\{T^{-1}(e_n)\}$ no es acotada en $c_0$. Por lo tanto, $T^{-1}$ no es continua.\\
            
            Si $T^{-1}$ fuese continua, entonces $Y$ sería isomorfo a $c_0$, y por tanto $Y$ sería un espacio de Banach, luego $Y$ sería cerrado en $l_2$. Como $T^{-1}$ no es continua, no podemos concluir que nada de forma directa. Veamos que $Y$ no es cerrado en $l_2$. Como $Y$ es denso en $l_2$, entonces $Y$ es cerrado si y solo si $Y = l_2$. Sabemos que $Y\subset l_2$, pero veamos que $Y\neq l_2$. Consideremos la sucesión $y\in l_2$ dada por:
            \begin{equation*}
                y(n) = \frac{1}{n} \qquad \forall n\in \mathbb{N}
            \end{equation*}

            Entonces $y\in l_2$, pero supongamos que $y\in Y$. Entonces existe $x\in c_0$ tal que $T(x) = y$, es decir:
            \begin{equation*}
                \frac{x(n)}{n} = \frac{1}{n} \implies x(n) = 1 \qquad \forall n\in \mathbb{N}
            \end{equation*}

            Esto es una contradicción, ya que $x\in c_0$ y por tanto $\{x(n)\}$ debe tender a cero. Por lo tanto, $y\notin Y$, y por tanto $Y\neq l_2$. Como $Y$ es denso en $l_2$, se concluye que $Y$ no es cerrado en $l_2$.
        \end{enumerate}
    \end{ejercicio}

    \begin{ejercicio}[1.5 puntos]
        Sean $X$ un espacio de Banach, $Y$ un espacio normado y $\cc{F}$ un subconjunto de $L(X,Y)$. Prueba que las siguientes afirmaciones son equivalentes:
        \begin{enumerate}[label=\alph*)]
            \item $\cc{F}$ está acotado.
            \item Para cada $x\in X$ y cada $g\in Y^\ast$ el conjunto $\{g(T(x)) : T\in \cc{F}\}$ está acotado.
        \end{enumerate}
        \begin{description}
            \item[$\Longrightarrow$)] Supongamos que $\cc{F}$ está acotado. Entonces, existe $M>0$ tal que:
            \begin{equation*}
                \|T\| \leq M \qquad \forall T\in \cc{F}
            \end{equation*}

            Sea $x\in X$ y $g\in Y^\ast$. Entonces, para cada $T\in \cc{F}$, por la continuidad de $g$ y la definición de la norma del funcional, se tiene:
            \begin{equation*}
                |g(T(x))| \leq \|g\|\cdot \|T(x)\|
            \end{equation*}

            Como $T\in \cc{F}$ es continuo, se tiene que:
            \begin{equation*}
                \|T(x)\| \leq \|T\|\cdot \|x\|
            \end{equation*}

            Como $\|T\| \leq M$, se concluye que:
            \begin{equation*}
                |g(T(x))| \leq \|g\|\cdot M \cdot \|x\|
            \end{equation*}

            Como $M, x, g$ son fijos, el conjunto $\{g(T(x)) : T\in \cc{F}\}$ está acotado.
        \end{description}
    \end{ejercicio}

    \begin{ejercicio}[2 puntos]
        Responde a las siguientes cuestiones indicando en cada caso, según proceda, un resultado de teoría que justifique tu respuesta o un contraejemplo, o bien dando una prueba.
        \begin{enumerate}[label=\alph*)]
            \item \textbf{[0.5 puntos]} ¿Puede definirse alguna norma completa en el espacio vectorial $c_{00}$ de las sucesiones casi nulas?
            \item \textbf{[0.5 puntos]} Si $H$ es un espacio de Hilbert y $M$ es un subespacio de $H$ verificando que $M^\perp = \{0\}$, entonces $M = H$.
            \item \textbf{[0.5 puntos]} Sea $\{f_n\}$ una sucesión puntualmente acotada de funciones lineales continuos en un espacio de Banach $X$, definiendo $T(x) = \{f_n(x)\}$ se obtiene un operador lineal continuo de $X$ en $t_\infty$.
            \item \textbf{[0.5 puntos]} Sean $\|\cdot \|$ y $|\|\cdot |\|$ dos normas completas no equivalentes en un espacio vectorial $X$. Entonces la aplicación identidad $I_X:(X,\|\cdot \|)\to (X,|\|\cdot |\|)$ no es continua.
        \end{enumerate}
    \end{ejercicio}

    \begin{ejercicio}[3 puntos]
        Escribe sobre alguno de los siguientes temas:
        \begin{enumerate}[label=\alph*)]
            \item Ortogonalidad. Teorema de la proyección ortogonal. Teorema de Riesz-Fréchet.
            \item Separación de conjuntos convexos en espacios normados.
            \item Lema de categoría de Baire. Teorema de Banach-Steinhaus.
        \end{enumerate}
    \end{ejercicio}

\end{document}
